\section{Descriptive Behavior of the Oracles}
\label{sec:results-oracle-behavior}

Table~\ref{tab:oracle-summary} summarizes the partitions produced for the joint configurations. Since both evaluated models use GPT-OSS 120B as the auxiliary model $\mathcal{A}$, their independently generated partitions are pooled. The counterfactual multiplier is $\sum_{s\in\mathbf{S}}(2^{|s|}-1)/|\mathbf{C}|$, relative to one atomic counterfactual per concept. Counterfactual counts refer to successfully written files, whereas the multiplier reflects the complete oracle partitions and is therefore not reduced by later generation failures. \emph{Concepts grouped} denotes the percentage of concepts assigned to non-singleton groups. All values are aggregated across the two runs with different evaluation models $\mathcal{M}$.

\begin{table}[htb]
  \centering
  \caption[Oracle partitions for the joint configurations]{Oracle partitions for the joint configurations, pooled across both evaluated models.}
  \label{tab:oracle-summary}
  \small
  \setlength{\tabcolsep}{6pt}
  \renewcommand{\arraystretch}{1.1}
  \rowcolors{2}{gray!8}{white}
  \begin{tabular}{@{}lrrr@{}}
    \hline
    \textbf{Measure}           & \textbf{BBQ}    & \textbf{MedQA}    & \textbf{New German Credit} \\
    \hline
    Group-size bound $n_{\max}$ & $4$             & $4$               & $3$                        \\
    Valid outputs              & $60/60$         & $58/60$           & $60/60$                    \\
    Mean group size            & $1.34$          & $1.45$            & $2.40$                     \\
    Groups of size $1/2/3$     & $137/66/2$      & $409/38/23$       & $60/60/180$                \\
    Groups of size $4/5/6$     & $0/0/0$         & $46/1/1$          & --                         \\
    Singleton groups           & $66.8\%$        & $79.0\%$          & $20.0\%$                   \\
    Concepts grouped           & $50.2\%$        & $45.4\%$          & $91.7\%$                   \\
    Counterfactuals            & $342$           & $1{,}468$         & $1{,}500$                  \\
    Counterfactuals per output & $5.70$          & $25.31$           & $25.00$                    \\
    Increase over atomic       & $1.27\times$    & $1.96\times$      & $2.08\times$               \\
    \hline
  \end{tabular}
\end{table}

The natural-language results contain many more singleton groups than the tabular partition, although roughly half of the concepts in both BBQ and MedQA still occur in non-singleton groups. Two anomalous MedQA groups contained five and six concepts, respectively, exceeding the imposed value of $n_{\max}=4$. BBQ's oracle partitions imply 349 counterfactuals, of which 342 were parsed and written successfully.

For the co-grouping heat maps, an OpenAI Codex large language model mapped the 46 raw BBQ and 236 raw MedQA labels to 10 and 19 high-level categories. It used only the labels, with one fixed mapping across all runs, and did not see group membership or calculated rates. Appendix~\ref{app:oracle-category-normalization} documents the method and definitions. This aggregation makes the resulting heat maps easier to interpret. Each cell is the percentage of within-question concept pairs placed in the same group, pooling both models. Gray denotes that no pair from the two categories co-occurred in an evaluated question. The use of conditional percentages prevents frequent categories from dominating solely because they occur more often.

\begin{figure}[htb]
  \centering
  \begin{adjustbox}{max width=\textwidth}
    \OracleHeatmapMedQA
  \end{adjustbox}
  \caption[Oracle co-grouping rates for MedQA]{Co-grouping rates of the MedQA oracle on LLM-normalized categories, pooled across both evaluated models.}
  \label{fig:oracle-heatmap-medqa}
\end{figure}

Figure~\ref{fig:oracle-heatmap-medqa} shows that the MedQA oracle groups concepts mainly within a clinical category. Pairs from the same category are placed in a common group in $22.5\%$ of eligible cases, compared with $6.0\%$ for pairs from different categories. The highest rates occur on the diagonal for \textit{Neurologic function} (13 of 16 eligible pairs) and \textit{Symptoms} (35 of 89), whereas \textit{Care context} and \textit{Demographics} are never grouped together despite 101 eligible pairs. In BBQ, only 14 of the 55 category combinations rest on at least ten eligible pairs, so most rates are based on few observations. Among these, \textit{Demographics} and \textit{Topics \& communication} are co-grouped most often, in 18 of 42 eligible pairs ($43\%$). The BBQ heat map is shown in Figure~\ref{fig:oracle-heatmap-bbq} in Appendix~\ref{app:oracle-heatmap-bbq}.

Figure~\ref{fig:german-credit-grouping-final} shows the deterministic tabular oracle's partition, reused for every New German Credit item: three groups of three concepts, one pair, and one singleton. This partition represents the final state of the grouping algorithm in Algorithm~\ref{alg:greedy_merge}. Node colors identify the five final groups, thick green edges are the accepted merge edges, and thin gray edges show other statistically retained pairwise associations. Edge labels report Cram\'er's $V$. The initialization and all seven accepted merge steps are documented in Appendix~\ref{app:oracle-partition-details}, and the corresponding retained-edge strengths are reported in Table~\ref{tab:credit-associations}.

\begin{figure}[htb]
  \centering
  \includestandalone[mode=buildnew,width=\textwidth,height=0.48\textheight,keepaspectratio]{graphs/german_credit_grouping}
  \caption{Final state of the constrained greedy grouping algorithm for New German Credit at $n_{\max}=3$.}
  \label{fig:german-credit-grouping-final}
\end{figure}
